% ===== PRÉAMBULE CANONIQUE — à coller VERBATIM en tête de CHAQUE .tex =====
\documentclass[11pt,a4paper]{article}
\usepackage[utf8]{inputenc}\usepackage[T1]{fontenc}\usepackage{lmodern}
\usepackage[french]{babel}
\usepackage{amsmath,amssymb,mathtools}\usepackage{icomma}
\usepackage[version=4]{mhchem}          % \ce{...} : équations & formules
\usepackage{siunitx}\sisetup{locale=FR} % \qty{}{}, \si{}, \ang{}
\usepackage{chemfig}                    % \chemfig{...} : structures développées
\usepackage{xcolor}\usepackage{colortbl}
\usepackage{tikz}\usepackage{pgfplots}\pgfplotsset{compat=1.18}
\usepackage[most]{tcolorbox}\usepackage{enumitem}\usepackage{array,booktabs}
\usepackage[a4paper,margin=2.2cm]{geometry}
\usetikzlibrary{arrows.meta,calc,angles,quotes,positioning,patterns,backgrounds,shapes.geometric}
\definecolor{primary}{RGB}{222,49,99}\definecolor{primarydark}{RGB}{150,22,55}
\definecolor{defcol}{RGB}{0,114,178}\definecolor{methcol}{RGB}{52,92,148}
\definecolor{excol}{RGB}{0,135,90}\definecolor{attncol}{RGB}{200,40,55}
\tcbset{boxbase/.style={enhanced,breakable,boxrule=0.9pt,arc=2.5pt,
  left=3mm,right=3mm,top=2.3mm,bottom=2.3mm,fonttitle=\bfseries,coltitle=white}}
% --- boîtes TUTORIEL (noms EXACTS, ne pas renommer) ---
\newtcolorbox{cle}[1][]{boxbase,colback=defcol!6,colframe=defcol,
  title={\textsc{À retenir}\ifx\relax#1\relax\relax\else\textmd{ : \itshape #1}\fi}}
\newtcolorbox{methode}[1][]{boxbase,colback=methcol!7,colframe=methcol,sharp corners,
  title={\textsc{Méthode}\ifx\relax#1\relax\relax\else\textmd{ --- \itshape #1}\fi}}
\newtcolorbox{exemplebox}[1][]{boxbase,colback=excol!5,colframe=excol,
  title={\textsc{Exemple}\ifx\relax#1\relax\relax\else\textmd{ : \itshape #1}\fi}}
\newtcolorbox{piege}[1][]{boxbase,colback=attncol!6,colframe=attncol,
  title={\textsc{Piège}\ifx\relax#1\relax\relax\else\textmd{ : \itshape #1}\fi}}
\newtcolorbox{remarque}[1][]{boxbase,colback=black!4,colframe=black!45,
  title={\textsc{Remarque}\ifx\relax#1\relax\relax\else\textmd{ : \itshape #1}\fi}}
% --- boîtes EXERCICE / CORRIGÉ (noms EXACTS, auto-comptées) ---
\newtcolorbox[auto counter]{exo}[1][]{boxbase,colback=primary!4,colframe=primary,
  title={\textsc{Exercice}~\thetcbcounter\ifx\relax#1\relax\relax\else\textmd{ --- \itshape #1}\fi}}
\newtcolorbox[auto counter]{corrige}[1][]{boxbase,colback=excol!4,colframe=excol,
  title={\textsc{Corrigé}~\thetcbcounter\ifx\relax#1\relax\relax\else\textmd{ --- \itshape #1}\fi}}
\newcommand{\R}{\mathbb{R}}\newcommand{\N}{\mathbb{N}}\newcommand{\Z}{\mathbb{Z}}\newcommand{\ds}{\displaystyle}
\begin{document}
\sisetup{per-mode=symbol}
\begin{corrige}[Masse molaire de petites molécules]
\begin{enumerate}[label=\alph*)]
  \item \ce{CH4} : $12+4\times1=\qty{16}{\gram\per\mole}$.
  \item \ce{CH4O} : $12+4\times1+16=\qty{32}{\gram\per\mole}$.
  \item \ce{C2H6} : $2\times12+6\times1=\qty{30}{\gram\per\mole}$.
\end{enumerate}
\end{corrige}

\begin{corrige}[Molécules oxygénées courantes]
\begin{enumerate}[label=\alph*)]
  \item \ce{C2H6O} : $2\times12+6+16=\qty{46}{\gram\per\mole}$.
  \item \ce{C3H6O} : $3\times12+6+16=\qty{58}{\gram\per\mole}$.
  \item \ce{C2H4O2} : $2\times12+4+2\times16=\qty{60}{\gram\per\mole}$.
\end{enumerate}
\end{corrige}

\begin{corrige}[Masse molaire ou masse moléculaire relative ?]
\ce{CH2O} : $12+2\times1+16=30$.
\begin{enumerate}[label=\alph*)]
  \item $M=\qty{30}{\gram\per\mole}$ (avec unité).
  \item $M_r=30$ (nombre \emph{sans} unité).
\end{enumerate}
\end{corrige}

\begin{corrige}[Pourcentage massique d'un élément]
$M(\ce{CH4})=\qty{16}{\gram\per\mole}$.
\begin{enumerate}[label=\alph*)]
  \item $\%_{\ce{C}}=\dfrac{1\times12}{16}\times100=\textbf{75\,\%}$.
  \item $\%_{\ce{H}}=\dfrac{4\times1}{16}\times100=\textbf{25\,\%}$.
\end{enumerate}
La somme fait bien $100\,\%$.
\end{corrige}

\begin{corrige}[De la structure à la masse molaire]
\begin{enumerate}[label=\alph*)]
  \item \ce{CH3-CO-CH3} : $3$~C, $6$~H, $1$~O $\Rightarrow \ce{C3H6O}$.
  \item $M=3\times12+6\times1+16=\qty{58}{\gram\per\mole}$.
\end{enumerate}
\end{corrige}

\end{document}
